\section{Câu hỏi trong tựa đề}
\label{sec:ch10-cau-hoi-trong-tua-de}

\index{LIME!biến thể|(}%
\index{khảo sát tài liệu}%
Muốn biết lời giải thích của LIME đổi giữa hai lần chạy vì đâu thì trước hết
phải biết LIME để ngỏ những gì.
Mục~\ref{sec:ch03-cac-tham-so-tu-do} đã liệt kê những chỗ mà LIME để người dùng
tự quyết: biểu diễn $x'$, phép dựng lại $z$, phân phối sinh nhiễu, hàm khoảng
cách $D$, độ rộng $\sigma$ của \tn{hàm nhân}{kernel}, số mẫu $N$ và độ dài $K$.
Mỗi lựa chọn ấy đứng thay cho một mệnh đề về hành vi của $f$, và không lựa chọn
nào tự chứng minh mệnh đề của nó.

Lĩnh vực đã phản ứng với danh sách đó, nhưng không phản ứng bằng cách kiểm
chứng từng lựa chọn, mà bằng cách đề xuất phương pháp mới. Trong chín năm kể
từ bài báo gốc, số công trình sửa một chỗ nào đó trong LIME đã đủ nhiều để chính
việc theo dõi chúng trở thành một vấn đề, và khảo sát của Knab và cộng sự ra đời
để giải quyết đúng vấn đề ấy~\autocite{p22whichlime}.

Quy mô của khối tài liệu được rà là phần đáng ghi lại. Khảo sát đi theo quy trình rà tài
liệu có cấu trúc của Webster và cộng sự, tìm quanh bài báo gốc bằng các cụm khóa
\enquote{LIME issues}, \enquote{LIME improvements} và \enquote{LIME
advancements}; vòng rà đầu tiên phủ các bài công bố trong khoảng từ 2016 tới
2025, rồi được mở rộng bằng cách lần theo tài liệu tham khảo của chính các bài
ấy. Các bản tiền ấn phẩm trên arXiv cũng được nhận nếu chúng có đóng góp mới và
liên quan~\autocite{p22whichlime}. Riêng bài báo gốc, tính tới ngày 21 tháng 1
năm 2025, đã được dẫn trong 21\,343 bài trên arXiv, và các tác giả ghi nhận rằng
họ có thể đã bỏ sót công trình vì phải rà một khối hơn 20\,000
tài liệu~\autocite{p22whichlime}.

Con số 21\,343 không đo chất lượng của LIME, và khảo sát cũng không dùng nó theo
nghĩa đó. Nó đo mức phổ biến, tức đo có bao nhiêu người đang đứng trên một
phương pháp mà chương~\ref{ch:lime} đã cho thấy là chưa được xác nhận. Đó là lý
do chương này đọc một khảo sát chứ không đọc thêm một phương pháp: cái cần nhìn
ở đây là hình dạng của cả nhánh nghiên cứu, không phải cơ chế của một phương
pháp trong nhánh ấy.

Khảo sát không đo gì cả, và điều đó ràng buộc mọi câu sau đây. Nó không dựng lại một biến thể nào,
không chạy một phép so nào giữa hai biến thể, và không có mục kết quả. Ba bảng
của nó phân loại công trình đã công bố: bảng thứ nhất đặt mỗi kỹ thuật vào một
nhóm vấn đề và một bước của quy trình, bảng thứ hai đặt mỗi kỹ thuật vào một
kiểu dữ liệu và một lĩnh vực ứng dụng, bảng thứ ba liệt kê các tính chất dùng
để đánh giá, lấy cấu trúc từ một khung có sẵn~\autocite{p22whichlime}. Vì vậy,
chỗ nào dưới đây nói LIME không ổn định thì đó là điều khảo sát ghi nhận từ tài
liệu nó dẫn, chứ không phải điều khảo sát đo được.
